% ============================================================
%  part2/ch08-fields-flows.tex
%  Chapter 8: Fields, Flows, and Conservation Laws
% ============================================================

\chapter{Fields, Flows, and Conservation Laws}
\label{ch:fields-flows}

\epigraph{%
  Nature does not make jumps.%
}{Gottfried Wilhelm Leibniz}

\begin{WhyMatters}
  Part~I showed that observation is projection.
  Part~II asks: what mathematical structure governs the
  object being projected?
  This chapter introduces the language of fields and flows
  that will carry through the rest of the book.
  Every RSVP field, every admissibility sheaf section, and
  every reconstruction operator will be a field in this sense.
\end{WhyMatters}

\section{From states to fields}

A state assigns a value to a system at a moment.
A field assigns a value to every point of a domain at every
moment.

\begin{definition}[Scalar field]
  A \emph{scalar field} on a domain $\Omega \subseteq \mathbb{R}^d$
  is a function $\phi : \Omega \times \mathbb{R} \to \mathbb{R}$.
\end{definition}

\begin{definition}[Vector field]
  A \emph{vector field} on $\Omega$ is a function
  $\mathbf{v} : \Omega \times \mathbb{R} \to \mathbb{R}^d$.
\end{definition}

\section{Conservation laws}

\begin{definition}[Continuity equation]
  A scalar field $\phi$ satisfies a \emph{conservation law} when
  \[
    \partial_t \phi + \nabla \cdot \mathbf{J} = \sigma,
  \]
  where $\mathbf{J}$ is the flux of $\phi$ and $\sigma$ is a
  source or sink term.
  When $\sigma = 0$, $\phi$ is \emph{strictly conserved}.
\end{definition}

\begin{proposition}[Global conservation]
  If $\phi$ satisfies the continuity equation with $\sigma = 0$
  and $\mathbf{J} \cdot \hat{n} = 0$ on $\partial\Omega$, then
  \[
    \frac{d}{dt} \int_\Omega \phi \, dx = 0.
  \]
\end{proposition}

\begin{proof}
  Integrate the continuity equation over $\Omega$ and apply the
  divergence theorem.
\end{proof}

\begin{ForwardRef}
  The RSVP entropy field $S(x,t)$ satisfies a modified
  continuity equation with a source term $\sigma(\Phi,\mathbf{v},L,D)$
  that is nonnegative when $S = 0$, ensuring $S \geq 0$ is
  preserved.
  The lamphron field $L = \ell^2$ satisfies a conservation
  law modified by advection.
  Both enter the free-energy functional in Part~V\@.
\end{ForwardRef}

\section{Gradient flows}

The most important class of evolution equations for this book
are gradient flows: dynamics that descend a functional.

\begin{definition}[Gradient flow]
  Given a functional $\mathcal{F}[\phi]$, the
  \emph{gradient flow} is
  \[
    \partial_t \phi = -M \frac{\delta \mathcal{F}}{\delta \phi},
  \]
  where $M > 0$ is a mobility coefficient.
\end{definition}

\begin{proposition}[Gradient flow dissipation]
  \label{prop:gf-dissipation}
  Along gradient flow solutions,
  \[
    \frac{d\mathcal{F}}{dt}
    = -M \int \left(\frac{\delta\mathcal{F}}{\delta\phi}\right)^2 dx
    \leq 0.
  \]
\end{proposition}

\begin{proof}
  $\frac{d\mathcal{F}}{dt}
  = \int \frac{\delta\mathcal{F}}{\delta\phi} \partial_t\phi\, dx
  = -M \int \left(\frac{\delta\mathcal{F}}{\delta\phi}\right)^2 dx
  \leq 0$.
\end{proof}

\begin{StatusBox}{Proven Framework}
  Proposition~\ref{prop:gf-dissipation} holds for any smooth
  functional $\mathcal{F}$ and any positive mobility $M$.
  It is the template for all energy-dissipation results in
  this book, including the RSVP free-energy dissipation
  theorem in Chapter~\ref{ch:energy-dissipation}.
\end{StatusBox}

\WhatSurvived{Field projection $\phi \mapsto \phi|_U$}{%
  Field values on $U$\\
  Local gradient information\\
  Conservation on $U$%
}{%
  Field values outside $U$\\
  Long-range correlations\\
  Boundary fluxes%
}{%
  Partial; requires boundary\\
  data and field equations.%
}{%
  $\ker(\cdot|_U)$: all field\\
  configurations supported\\
  outside $U$.%
}
